\documentclass[9pt]{extarticle}
\usepackage[utf8]{inputenc}
\usepackage[T1]{fontenc}
\usepackage[spanish]{babel}
\usepackage[a4paper,landscape,margin=1.1cm,top=1.2cm,bottom=1.3cm]{geometry}
\usepackage{lmodern}
\usepackage{array}
\usepackage{longtable}
\usepackage{xcolor}
\usepackage{colortbl}
\usepackage{fancyhdr}

\definecolor{gris}{gray}{0.88}
\definecolor{grisclaro}{gray}{0.95}
\setlength{\parindent}{0pt}
\renewcommand{\arraystretch}{1.25}
\setlength{\tabcolsep}{4pt}

\pagestyle{fancy}
\fancyhf{}
\renewcommand{\headrulewidth}{0.6pt}
\setlength{\headheight}{24pt}
\lhead{\textbf{COLEGIO HISPANO AMERICANO RÍO VIEJO}\\ \small\itshape Al Servicio de Dios y la Educación}
\rhead{\small Coordinación UTP}
\cfoot{\footnotesize\itshape “Antes sed los unos con los otros benignos, misericordiosos, perdonándoos los unos a los otros
como también Dios os perdonó en Cristo” — Efesios 4:32}

\newcolumntype{C}[1]{>{\raggedright\arraybackslash}p{#1}}
\newcommand{\oa}[1]{\rowcolor{grisclaro}\multicolumn{6}{|l|}{\textbf{#1}} \\ \hline}
\newcommand{\pond}[1]{\rowcolor{gris}\multicolumn{6}{|l|}{\textbf{Ponderación #1}} \\ \hline}
\newcommand{\crit}[2]{\textbf{#1}\newline #2}

\begin{document}

\begin{center}
  {\large\bfseries RÚBRICA UNIDAD N°4 LENGUAJE}\\[2pt]
  {\bfseries “EL CABALLERO DE LA ARMADURA OXIDADA” ROBERT FISHER}\\[2pt]
  Profesora: Elizabeth Viscay / Educ. Dif. Valeska Cáceres
\end{center}

\vspace{2pt}
\begin{tabular}{|p{12.8cm}|p{6cm}|p{6.2cm}|}
\hline
\textbf{NOMBRE:} & \textbf{CURSO:} & \textbf{FECHA:} \\[8pt] \hline
\end{tabular}

\vspace{4pt}
\textbf{Objetivo/s:} Expresarse de manera clara y efectiva en exposiciones orales para comunicar temas de su interés.

\vspace{4pt}
\begin{longtable}{|C{4.9cm}|C{4.45cm}|C{4.45cm}|C{4.45cm}|C{4.45cm}|>{\centering\arraybackslash}p{1.55cm}|}
\hline
\rowcolor{gris}
\textbf{Criterio o dimensión} & \textbf{EXCELENTE (3)} & \textbf{SATISFACTORIO (2)} & \textbf{EN PROCESO DE LOGRO (1)} & \textbf{INSATISFACTORIO (0)} & {\footnotesize\textbf{PUNTAJE}} \\ \hline
\endhead
& El/la estudiante demuestra un dominio total y autónomo del indicador. Cumple de manera sobresaliente, precisa y detallada con todo lo solicitado, superando las expectativas básicas y sin cometer errores.
& El/la estudiante cumple de manera adecuada y correcta con lo solicitado en el indicador. Muestra un buen dominio, aunque puede presentar omisiones menores o detalles mínimos por profundizar que no afectan el resultado general.
& El/la estudiante cumple de forma parcial, incompleta o inconsistente con el indicador. Demuestra un dominio débil o inicial, requiriendo apoyo o guía para alcanzar el estándar esperado.
& El/la estudiante no cumple con lo solicitado en el indicador, presenta información errónea o simplemente omite/no realiza o no se observa el aspecto evaluado.
& \\ \hline

\pond{80\%}
\oa{OA 4 Analizar aspectos relevantes de las narraciones leídas para profundizar su comprensión.}

\crit{Descripción física y psicológica de los personajes}{Describe física y psicológicamente a los personajes de acuerdo con las actitudes que toman en el relato y sus reacciones frente a los problemas. (Mínimo 3 características).}
& Describe detalladamente al menos 3 características (físicas y psicológicas) de los personajes, vinculándolas claramente con las actitudes que toman en el relato y sus reacciones frente a los problemas.
& Menciona las 3 características solicitadas (físicas y/o psicológicas), pero explica de forma parcial su relación con las actitudes o problemas del relato.
& Menciona menos de 3 características de los personajes o solo nombra rasgos superficiales sin vincularlos a las acciones del libro.
& No describe a los personajes o la información presentada es errónea e irrelevante.
& \\ \hline

\crit{Opinión fundamentada del libro}{Elabora una opinión fundamentada del libro a partir de la información leída del texto, sus experiencias y conocimientos.}
& Elabora una opinión personal clara y sólidamente fundamentada a partir de la información del texto, sus propias experiencias y conocimientos previos.
& Emite una opinión sobre el libro fundamentada principalmente en el texto, con escasa vinculación a sus experiencias o conocimientos personales.
& Da una opinión muy breve o vaga (ej. “me gustó” / “no me gustó”) con argumentos débiles o sin fundamentación suficiente.
& No manifiesta una opinión sobre la obra.
& \\ \hline

\crit{Argumentación y fundamentación con ejemplos del texto}{Argumenta su opinión o postura frente a las acciones o actitud del personaje principal y la fundamentan con ejemplos del texto.}
& Argumenta con claridad su postura frente a las acciones o actitudes del personaje principal, respaldando sus ideas con ejemplos precisos extraídos del libro.
& Argumenta su postura frente al personaje principal, pero los ejemplos del texto son generales o incompletos.
& Expresa una postura pero no logra fundamentarla o los ejemplos citados no corresponden al texto.
& No argumenta su postura ni menciona ejemplos del texto.
& \\ \hline

\crit{Propuesta de solución al conflicto/problema del personaje}{Cómo solucionarías el problema del personaje principal.}
& Propone una solución coherente, creativa y argumentada para resolver el conflicto del personaje principal, demostrando comprensión profunda de la trama.
& Propone una solución adecuada al problema del personaje principal, pero con una justificación breve o simple.
& Menciona una solución poco clara, poco coherente con la historia o poco aplicable al conflicto del personaje.
& No plantea ninguna solución al problema del personaje.
& \\ \hline

\oa{OA 9 Desarrollar el gusto por la lectura, leyendo habitualmente diversos textos.}

\crit{Expresión de apreciación personal}{Expresan qué fue lo que más les gustó de la historia y lo que menos, con fundamentos y ejemplos del texto.}
& Expresa de forma clara qué fue lo que más y lo que menos le gustó de la historia, fundamentando ambas posturas con ejemplos concretos del libro.
& Expresa lo que más y lo que menos le gustó de la historia, pero sólo fundamenta una de sus apreciaciones con ejemplos.
& Menciona qué le gustó o disgustó sin dar fundamentos ni ejemplos de la obra.
& No señala aspectos positivos ni negativos de la lectura.
& \\ \hline

\oa{OA 28 Expresarse de manera clara y efectiva en exposiciones orales para comunicar temas de su interés.}

\crit{Lenguaje no verbal}{Adoptan gestos y posturas adecuadas a la situación, hacen contacto visual con la audiencia, mantiene buena postura y uso del espacio en la sala.}
& Adopta posturas y gestos acordes a la situación, mantiene contacto visual constante con la audiencia y utiliza adecuadamente el espacio de la sala.
& Mantiene buena postura y contacto visual intermitente con la audiencia; su uso del espacio es limitado pero adecuado.
& Muestra rigidez o movimientos distractores, escaso contacto visual y permanece estático durante la exposición.
& Su lenguaje no verbal es inadecuado (da la espalda al público, se recuesta o distrae constantemente a la audiencia).
& \\ \hline

\crit{Lenguaje verbal}{Emplean un volumen, modulación, pronunciación y entonación adecuados.}
& Emplea de manera constante un volumen audible, modulación clara, pronunciación correcta y entonación variada acorde al discurso.
& Su volumen y pronunciación son adecuados la mayor parte del tiempo, aunque presenta leves fallas en modulación o entonación.
& Presenta volumen bajo, titubeos o problemas recurrentes de modulación que dificultan el entendimiento.
& El volumen es inaudible o la pronunciación impide comprender la exposición.
& \\ \hline

\crit{Uso de material de apoyo}{Utilizan material de apoyo que complemente la exposición (video, ppt, cartulina, etc.)}
& Utiliza un recurso (PPT, papelógrafo/cartulina, video) limpio, legible y atractivo, que complementa y enriquece de manera directa la exposición.
& Presenta material de apoyo relacionado con la exposición, pero lo utiliza poco durante el desarrollo de la misma.
& El material de apoyo es escaso, desordenado o apenas aporta a la exposición oral.
& No presenta material de apoyo.
& \\ \hline

\crit{Fluidez y autonomía}{Expresan las ideas e información sin leer.}
& Expone de forma fluida y natural sus ideas, sin recurrir a la lectura de fichas, diapositivas o apuntes.
& Expone principalmente con sus palabras, recurriendo a leer apuntes o diapositivas sólo de manera puntual.
& Depende constantemente de la lectura del material de apoyo para poder realizar su exposición.
& Se limita a leer textualmente todo su material o ficha durante la presentación.
& \\ \hline

\crit{Vocabulario y registro de habla}{Utilizan vocabulario variado, preciso y un registro formal adecuado a la situación comunicativa.}
& Utiliza un vocabulario variado, preciso y mantiene un registro formal adecuado a la situación académica.
& Utiliza un vocabulario adecuado y registro formal, con escasas repeticiones de palabras o muletillas.
& Emplea vocabulario limitado, informal o hace uso recurrente de muletillas durante su presentación.
& Emplea un registro informal e inadecuado para una exposición escolar.
& \\ \hline

\crit{Estructura de la exposición}{Organizan su presentación con una introducción, desarrollo y cierre.}
& Estructura la presentación diferenciando claramente la introducción (saludo/presentación), el desarrollo del tema y un cierre/conclusión.
& Presenta introducción, desarrollo y cierre, pero la transición entre las partes resulta poco clara.
& Omite la introducción o el cierre de su exposición, pasando directamente a los contenidos o finalizando bruscamente.
& La presentación no posee una estructura reconocible.
& \\ \hline

\crit{Pertinencia de la información presentada}{Realiza una exposición oral en la que se presenta información que se relaciona con el tema.}
& Toda la información expuesta está directamente centrada y relacionada con la temática de la obra y los requerimientos del trabajo.
& La mayor parte de la información se relaciona con el tema, incluyendo algunos datos secundarios innecesarios.
& Incluye información poco relevante o desviada del tema central de la exposición.
& La información presentada no guarda relación con la obra asignada.
& \\ \hline

\crit{Dominio del tema}{Demuestra dominio del tema, evidenciando conocimiento de conceptos e información presentada.}
& Demuestra amplio dominio del contenido, explicando conceptos y detalles clave del libro con seguridad y soltura.
& Demuestra conocimiento del libro, aunque muestra dudas en algunos conceptos o momentos puntuales.
& Muestra vacíos evidentes de lectura o confusión en conceptos e información relevante.
& Se evidencia que no leyó el libro ni maneja los contenidos mínimos.
& \\ \hline

\multicolumn{5}{|r|}{\textbf{Puntaje 80\%}} & \\ \hline
\multicolumn{5}{|r|}{\textbf{Nota 80\%}} & \\ \hline

\pond{20\%}

\crit{Presentación}{Saludan, indican sus nombres, el nombre del libro y autor.}
& Saluda formalmente, indica claramente su nombre completo, el título exacto del libro y el autor.
& Saluda e indica su nombre y el título del libro, pero omite señalar al autor.
& Omite datos clave en la presentación inicial (como su nombre o el título).
& Inicia la exposición sin saludar ni presentar la obra.
& \\ \hline

\crit{Presentación personal}{Presentación personal acorde a la situación (uniforme y aseo personal).}
& Se presenta con uniforme completo ordenado y óptima higiene/aseo personal, de acuerdo al reglamento escolar.
& Se presenta con el uniforme del colegio, pero con detalles de orden menor.
& Asiste con prendas que no corresponden al uniforme escolar o de manera desordenada.
& No cumple con las normas básicas de presentación personal del establecimiento.
& \\ \hline

\crit{Respeto hacia los compañeros/as}{Muestran respeto por las exposiciones de sus compañeros.}
& Mantiene una actitud de escucha activa y silencio respetuoso durante las exposiciones de sus pares.
& Escucha a sus compañeros/as, requiriendo un llamado de atención leve.
& Genera ruidos, conversa o distrae esporádicamente durante la presentación de otros.
& Interrumpe constantemente o muestra falta de respeto hacia las exposiciones ajenas.
& \\ \hline

\crit{Ortografía en el material de apoyo}{El material de apoyo escrito presenta ortografía puntual, literal y acentual correcta.}
& El material visual impreso o proyectado no presenta errores de ortografía puntual, literal y acentual.
& El material presenta entre 1 y 3 errores ortográficos en total.
& El material presenta entre 4 y 6 errores ortográficos.
& El material presenta más de 6 errores ortográficos o carece de revisión.
& \\ \hline

\crit{Manejo del tiempo y adaptabilidad}{Capacidad de respetar el tiempo asignado para su exposición y mantener la calma o adaptarse si surge algún inconveniente técnico/material.}
& Respeta el tiempo asignado para la exposición y reacciona con calma y flexibilidad ante cualquier imprevisto técnico o material.
& Excede o reduce levemente el tiempo asignado, pero se adapta de buena forma a imprevistos.
& No gestiona bien su tiempo (muy breve o demasiado extenso) o se muestra bloqueado/a ante fallas técnicas.
& Abandona la exposición por falta de tiempo o incapacidad de manejar un inconveniente menor.
& \\ \hline

\crit{Atención a preguntas e interacción}{Disponibilidad y disposición para recibir preguntas o comentarios de sus compañeros/as al finalizar la presentación y responder de forma respetuosa y receptiva.}
& Muestra excelente disposición para recibir preguntas/comentarios al finalizar y responde con respeto, fundamentación y receptividad.
& Acepta responder preguntas de sus compañeros/as y entrega respuestas breves pero pertinentes.
& Responde con inseguridad, de forma muy escueta o muestra incomodidad al recibir preguntas.
& Se niega a recibir o responder preguntas de la audiencia.
& \\ \hline

\crit{Responsabilidad y puntualidad}{Cumplimiento de los plazos estipulados para presentar su trabajo y tener listos los materiales de apoyo.}
& Presenta su trabajo en la fecha y hora estipuladas, teniendo todos sus materiales listos con anticipación.
& Presenta el día estipulado, pero sufre retrasos organizando sus materiales al momento de iniciar.
& Presenta en la fecha establecida, pero olvida parte de sus materiales de apoyo.
& No presenta en la fecha correspondiente sin justificación previa o médica.
& \\ \hline

\crit{Comprensión del contenido expuesto (Uso de fuentes)}{Se evidencia uso de conceptos o frases sacadas de internet que no comprende.}
& Utiliza sus propias palabras y muestra comprensión cabal de todas las frases, palabras y conceptos expuestos.
& Utiliza un lenguaje propio en su mayoría, con alguna frase técnica o extraída que logra explicar al ser consultado.
& Utiliza conceptos o frases complejas copiadas de internet que no logra explicar correctamente.
& Evidencia copiado textual masivo desde internet sin entender lo que está exponiendo.
& \\ \hline

\multicolumn{5}{|r|}{\textbf{Puntaje 20\%}} & \\ \hline
\multicolumn{5}{|r|}{\textbf{Nota 20\%}} & \\ \hline
\multicolumn{5}{|r|}{\textbf{NOTA FINAL}} & \\ \hline
\end{longtable}

\end{document}
